\chapter{Improving GRPO: Diagnostics, Stability and Its Descendants}\label{ch:grpo2}

\Cref{ch:grpo} ended with a working GRPO loop and a warning: plain GRPO is fine for a few
dozen steps but can become unstable over hundreds. This chapter is about the difference
between an algorithm that works in a notebook and one you can run for a week. Three skills
are involved. The first is \emph{reading} a run: which metrics to log, what healthy and
unhealthy curves look like, and why the loss is the least informative number on the
dashboard. The second is \emph{stabilising} it with the two mechanisms of the original GRPO
objective we skipped so far, clipped importance ratios and a KL penalty, and with auxiliary
rewards. The third is knowing the \emph{design space} that has opened up since DeepSeek-R1:
DAPO, Dr.~GRPO, GSPO, CISPO, GDPO and the DeepSeek-V3.2 corrections each change one
specific term of \cref{eq:grpo:full}, and each has a cost that a systems person should be
able to estimate.

Every experiment in this chapter is a run of \code{labs/lab26\_grpo\_plus/run.py} on the toy
addition task of \cref{ch:grpo}, from the same SFT-warm-started model, with two seeds per
configuration, measured on a CPU in this book's build environment. The toy cannot reproduce
every failure of a 0.6B model on MATH. When it fails to, we say so and explain what the real
run does and why. As with any RL result, single runs are noisy: treat differences smaller
than the seed-to-seed spread as noise.

\begin{figure}[t]
\centering
\resizebox{\linewidth}{!}{%
\begin{tikzpicture}[t/.style={draw,rounded corners=3pt,minimum width=34mm,minimum height=12mm,
  align=center,font=\small},num/.style={circle,fill=#1,text=white,font=\scriptsize\bfseries,inner sep=1.4pt},
  >={Stealth[length=2mm]}]
\node[font=\small\bfseries,keygreen] at (0,1.25) {observe};
\node[font=\small\bfseries,sysblue] at (4.2,1.25) {constrain};
\node[font=\small\bfseries,napkin] at (8.4,1.25) {shape \& extend};
\node[t,fill=keygreen!10,draw=keygreen] (a) at (0,0) {basic metrics\\{\scriptsize reward, length, eval}};
\node[t,fill=keygreen!10,draw=keygreen] (b) at (0,-1.7) {advanced metrics\\{\scriptsize advantages, entropy}};
\node[t,fill=sysblue!10,draw=sysblue] (c) at (4.2,0) {clipped ratios\\{\scriptsize trust region per step}};
\node[t,fill=sysblue!10,draw=sysblue] (d) at (4.2,-1.7) {KL to a reference\\{\scriptsize leash over the run}};
\node[t,fill=napkin!12,draw=napkin] (e) at (8.4,0) {format reward\\{\scriptsize \texttt{<think>} tags}};
\node[t,fill=napkin!12,draw=napkin] (f) at (8.4,-1.7) {GRPO descendants\\{\scriptsize DAPO, Dr.\,GRPO, GSPO, \dots}};
\foreach \n/\c/\i in {a/keygreen/1,b/keygreen/2,c/sysblue/3,d/sysblue/4,e/napkin/5,f/napkin/6}
  \node[num=\c] at (\n.north west) {\i};
\draw[->,thick,gray] (a)--(b); \draw[->,thick,gray] (b.east)--(c.west);
\draw[->,thick,gray] (c)--(d); \draw[->,thick,gray] (d.east)--(e.west); \draw[->,thick,gray] (e)--(f);
\node[draw=labpurple,dashed,rounded corners,fit=(a)(f),inner sep=7pt,
  label={[font=\scriptsize,labpurple]below:Chapter 37's GRPO loop is the starting point; Lab 38 ablates every box}] {};
\end{tikzpicture}}
\caption{Roadmap. We first learn to observe a run (1--2), then constrain how far each
update and the whole run may move the policy (3--4), and finally shape the reward and
survey the family of algorithms that grew out of GRPO (5--6).}
\label{fig:grpo2:roadmap}
\end{figure}

% ======================================================================
\section{From a working loop to a robust one}\label{sec:grpo2:intro}

It helps to name what we implemented in \cref{ch:grpo}. With summed log-probabilities, no
ratio, no clipping and no KL, the update is REINFORCE with a group-normalised baseline. Its
weaknesses are well known from classical RL. Nothing limits the size of one update except
the learning rate, so a single unlucky group with extreme advantages can move the policy a
long way. Nothing anchors the policy over the long run, so it can drift into degenerate text
once rewards stop carrying signal. And every rollout is used exactly once, which is wasteful
when rollouts are the expensive part. PPO's two devices, ratio clipping and a KL anchor,
address the first two weaknesses and make the third (reusing rollouts) safe. The rest of this
chapter adds them one at a time and watches what changes.

% ======================================================================
\section{Tracking GRPO performance metrics}\label{sec:grpo2:metrics}

\subsection{Executing a training run}\label{sec:grpo2:exec}

Long RL runs belong in scripts, not notebooks. A script survives a dropped SSH session
(run it under \code{tmux} or a job scheduler), writes structured logs that can be plotted
while the run is still going, and makes the configuration explicit. The lab's driver takes
the variant name and seeds on the command line and writes one CSV row per step:

\begin{lstlisting}[style=shell,caption={Running GRPO variants from the terminal.},label={lst:grpo2:run}]
cd code/labs/lab38_grpo_plus
python run.py --list                                  # show all variants
python run.py --variants baseline,clip_reuse --seeds 0,1 --steps 150 --out out18
\end{lstlisting}

The CSV columns are produced by \code{GRPOTrainer.step} and are the vocabulary of the rest
of the chapter: \code{loss}, \code{reward\_avg}, \code{correct\_avg}, \code{format\_avg},
\code{adv\_avg}, \code{adv\_std}, \code{entropy\_avg}, \code{resp\_len},
\code{zero\_var\_frac}, \code{policy\_ratio}, \code{clip\_frac}, \code{kl}, the per-phase
wall-clock times \code{t\_rollout}, \code{t\_logprob}, \code{t\_update}, and periodically
\code{eval\_acc}. Writing a row per step is cheap. Recomputing a missing metric afterwards
means re-running the experiment.

\begin{systemslens}[logging is part of the system]
At scale, metrics come from hundreds of ranks. Reduce them \emph{before} logging: all-reduce
the per-rank sums and counts of reward, tokens and entropy, never the per-rank means, which
are biased when ranks hold different numbers of tokens. Log asynchronously so a slow
metrics backend never stalls the trainer. Also log the \emph{throughput} metrics (tokens
generated per second, tokens trained per second, sampler idle fraction). An RL run that
learns correctly but uses its GPUs at 15\% is a bug too, just a more expensive one.
\end{systemslens}

\subsection{Inspecting the run}\label{sec:grpo2:inspect}

RL metrics are noisy step to step, so we always plot the raw series faintly and a trailing
moving average on top:

\begin{lstlisting}[caption={Trailing moving average and a multi-panel metrics plot (\texttt{llmsys/rl/metrics.py}).},label={lst:grpo2:plot}]
def moving_average(values, window):
    out, s = [], 0.0
    for i, v in enumerate(values):
        s += v
        if i >= window:
            s -= values[i - window]
        out.append(s / min(i + 1, window))   # shorter window at the start
    return out

plot_metrics("out18/baseline_s0.csv",
             ["reward_avg", "adv_std", "entropy_avg", "zero_var_frac"])
\end{lstlisting}

What should a healthy run look like? Four series carry most of the information:
\begin{itemize}
\item \textbf{Loss.} A sanity check, nothing more. GRPO's loss is a weighted
log-likelihood whose weights average to zero within every group. It is \emph{not} expected
to decrease, and its magnitude reflects how much groups disagreed, not how good the
model is. What \emph{is} alarming is a sudden order-of-magnitude jump, or frequent spikes
late in a run.
\item \textbf{Average reward.} Should trend upwards. A reward that sits at 1.0 is good news
with a catch: every group is then zero-variance (\cref{sec:grpo:adv}) and the run has stopped
learning. Stop it, or move to harder prompts (curriculum).
\item \textbf{Response length.} On real reasoning tasks it typically \emph{grows} as the
model learns to think longer. A length pinned at the token budget means rollouts are being
truncated, so they earn no reward and push the policy in confusing directions. A collapsing
length can mean the policy found a shortcut.
\item \textbf{Held-out accuracy.} The quantity we care about, measured rarely because it is
expensive (\cref{ch:eval}). A reward that keeps rising while held-out accuracy falls is the
signature of reward hacking or overfitting to the training prompts.
\end{itemize}
A typical long run of plain GRPO on Qwen3-0.6B with MATH looks like this: reward and held-out
accuracy both rise for the first few dozen steps. Over the following hundreds of steps the loss
begins to spike more often, response length first grows and then falls, and MATH-500
accuracy declines from its peak, even though it is measured on the same 500 problems
throughout. A decline on a fixed evaluation set cannot be explained by later training
prompts being harder. It points at the algorithm.

\begin{napkinmath}[the cost of watching]
Evaluating on MATH-500 with a 2048-token budget is 500 rollouts, about $10^6$ generated
tokens. For a 0.6B model in an engine at about \num{10000} output tokens/s that is around 100\,s,
cheap if done every 50 steps. With an in-loop, unbatched HF-style generator at about 30 tokens/s
it is roughly 9 hours, which is why naive \code{--eval\_every} flags slow training to a crawl.
Batch the evaluation (\cref{ch:serving}), run it on a separate replica from saved
checkpoints, or evaluate on a fixed random subset of 100 problems and accept wider error
bars ($\pm$10 points at 95\% confidence for accuracies near 50\%).
\end{napkinmath}

% ======================================================================
\section{Tracking more advanced metrics}\label{sec:grpo2:advanced}

Two quantities computed inside the GRPO step are worth logging even though they never
appear in the loss: summary statistics of the advantages, and the entropy of the policy's
next-token distributions.

\subsection{Advantage statistics}\label{sec:grpo2:advstats}

\begin{lstlisting}[caption={Advantage statistics.},label={lst:grpo2:advstats}]
def advantage_stats(rewards):
    r = torch.as_tensor(rewards, dtype=torch.float32)
    adv = (r - r.mean()) / (r.std() + 1e-4)
    return adv, adv.mean().item(), adv.std().item()
\end{lstlisting}

\emph{The mean is zero by construction.} Each group's advantages are mean-centred, so their
mean is zero up to floating-point error. Logging it is a free unit test: a drifting
\code{adv\_avg} means advantages are being normalised over the wrong axis (for example,
across the whole batch instead of within groups) or masked incorrectly.

\emph{The standard deviation measures the learning signal.} For a group with mixed rewards
the std of the normalised advantages is $\approx1$. For a zero-variance group it is 0.
Averaged over a batch, \code{adv\_std} therefore tracks roughly the fraction of groups that
carry signal. Values near 1 mean healthy, well-scaled updates; values falling towards 0 mean
the signal is vanishing (rewards all 0 or all 1); erratic spikes mean a few groups dominate.
In practice it is clearer to log the fraction directly, which our trainer does as
\code{zero\_var\_frac}. Read it together with the average reward: zero variance at reward 0
means ``too hard'', at reward 1 ``solved''.

\begin{figure}[ht]
\centering
\begin{tikzpicture}
\begin{axis}[width=0.7\linewidth,height=4.6cm,xlabel={per-prompt success rate $p$},
  ylabel={$P(\text{zero-variance group})$},xmin=0,xmax=1,ymin=0,ymax=1,grid=major,grid style={gray!15},
  legend style={font=\scriptsize,at={(0.5,0.97)},anchor=north},tick label style={font=\scriptsize}]
\addplot[very thick,warnred,domain=0:1,samples=101] {x^4+(1-x)^4}; \addlegendentry{$G=4$}
\addplot[very thick,sysblue,domain=0:1,samples=101] {x^8+(1-x)^8}; \addlegendentry{$G=8$}
\addplot[very thick,keygreen,domain=0:1,samples=101] {x^16+(1-x)^16}; \addlegendentry{$G=16$}
\end{axis}
\end{tikzpicture}
\caption{Probability that a group of $G$ binary-reward rollouts carries no gradient,
$p^G+(1-p)^G$. Very easy and very hard prompts waste their rollouts. Larger groups help in
the middle of the range but not at the ends, which is why prompt selection (curricula,
DAPO's dynamic sampling) matters as much as $G$.}
\label{fig:grpo2:zerovar}
\end{figure}

\subsection{Entropy}\label{sec:grpo2:entropy}

Entropy measures how spread out the policy's next-token distribution is,
\begin{equation}
H(\pi(\cdot\mid s))=-\sum_{v\in\mathcal V}\pi(v\mid s)\log\pi(v\mid s)\in[0,\log|\mathcal V|].
\end{equation}
It is 0 when one token has all the mass and $\log|\mathcal V|$ when the distribution is
uniform. RL needs \emph{some} entropy, since rollouts in a group only differ if the policy
is uncertain somewhere, and every policy-gradient update tends to reduce it, since it
concentrates mass on rewarded choices. Logging entropy shows which of the two forces is
winning.

\begin{figure}[ht]
\centering
\begin{tikzpicture}
\begin{axis}[ybar,width=0.75\linewidth,height=4.6cm,bar width=11pt,
  symbolic x coords={7,8,6,9,1},xtick=data,xlabel={candidate next token after \texttt{\textbackslash boxed\{}},
  ymin=0,ymax=0.7,legend style={font=\scriptsize,at={(0.98,0.97)},anchor=north east},
  nodes near coords,every node near coord/.append style={font=\tiny},tick label style={font=\scriptsize},
  enlarge x limits=0.15]
\addplot[fill=sysblue!55,draw=none] coordinates {(7,0.603) (8,0.222) (6,0.135) (9,0.030) (1,0.011)};
\addplot[fill=napkin!70,draw=none] coordinates {(7,0.305) (8,0.334) (6,0.270) (9,0.105) (1,0.050)};
\legend{$p_v$,$-p_v\log p_v$}
\end{axis}
\end{tikzpicture}
\caption{Entropy of one next-token distribution (a 5-token toy vocabulary, logits
$2,1,0.5,-1,-2$). The entropy is the sum of the orange bars, $H=1.06$ nats, out of a maximum
of $\log5=1.61$. Note that the most likely token does not contribute the most: the term
$-p\log p$ peaks at $p=1/e$.}
\label{fig:grpo2:entropybar}
\end{figure}

For orientation, with natural logarithms: per-token entropies near 0--0.5 mean the model is
almost deterministic at that step; 1--2 means a handful of plausible tokens, which is typical
of healthy text; values approaching $\log|\mathcal V|$ (11.9 for Qwen3) mean noise. Averaged
over a reasoning trace, most tokens are near-deterministic and a few ``forking'' tokens carry
most of the entropy. Two failure modes show up in the average:
\emph{entropy collapse} (the policy stops exploring, groups become identical, learning stalls)
and \emph{entropy explosion} (updates have damaged the model and it emits random tokens).

Since \code{log\_softmax} already gives all log-probabilities, entropy costs one extra
elementwise product and reduction. \Cref{lst:grpo2:entropy} extends the sequence log-prob
function of \cref{ch:grpo} to return both.

\begin{lstlisting}[caption={Sequence log-probability and mean per-step entropy in one pass.},label={lst:grpo2:entropy}]
def sequence_logprob_and_entropy(model, token_ids, prompt_len):
    logits = model(token_ids[None])[0].float()
    logp = torch.log_softmax(logits, dim=-1)
    picked = logp[:-1].gather(1, token_ids[1:, None]).squeeze(-1)[prompt_len-1:]
    steps = logp[:-1][prompt_len - 1:]           # (|y|, V): one row per answer step
    if steps.shape[0] == 0:                      # model emitted nothing
        return picked.sum(), picked.new_tensor(0.0)
    ent = -(steps.exp() * steps).sum(-1)         # H at every answer step
    return picked.sum(), ent.mean().detach()     # entropy is a metric, not a loss
\end{lstlisting}

\begin{systemslens}[entropy is a full-vocabulary reduction]
The log-probability of the sampled token needs one element per row; entropy needs \emph{all}
$|\mathcal V|$ of them. In a trainer that uses a fused ``linear + cross-entropy'' kernel to
avoid materialising logits, adding an entropy metric can silently bring back the
$(N\cdot T)\times|\mathcal V|$ float32 tensor: $8\times1024\times\num{151936}\times4$\,B
$\approx5$\,GB. Compute it in row chunks (\code{chunked\_logprob\_entropy} in the lab bounds
the temporary to \code{chunk}$\times|\mathcal V|$), under \code{no\_grad}, and only on the
``old log-prob'' pass that runs anyway, as \code{GRPOTrainer} does. Some trainers sample a
subset of positions for entropy, which is fine for a metric.
\end{systemslens}

\subsection{Plotting the additional metrics}\label{sec:grpo2:plotadv}

\Cref{fig:grpo2:baseline} shows the advanced metrics of our baseline run (\cref{ch:grpo}'s
algorithm, seed 0). The average advantage stays at zero throughout, as it must. The other
three curves tell one story. At the start about half the groups are zero-variance because
all eight rollouts are \emph{wrong} (mostly unboxed). Within twenty steps the format is
learned, the zero-variance fraction drops to about 15\%, and the advantage std peaks near
0.87. From then on more and more prompts are solved by \emph{all} eight rollouts, the
zero-variance fraction climbs back above 50\%, and the advantage std declines. By the end,
half of every step's rollouts produce no gradient. Entropy decays smoothly from 0.53 to 0.09
nats per token, a policy sharpening on a task with exactly one right answer per prompt. That
is healthy here, but it would be a warning sign on open-ended reasoning, where exploration
must survive.

\begin{figure}[ht]
\centering
\begin{tikzpicture}
\begin{groupplot}[group style={group size=2 by 2,horizontal sep=1.5cm,vertical sep=1.1cm},
  width=0.48\linewidth,height=4cm,xmin=0,xmax=150,grid=major,grid style={gray!15},
  tick label style={font=\scriptsize},label style={font=\small}]
\nextgroupplot[ylabel={advantage mean},ymin=-0.5,ymax=0.5]
\addplot[sysblue!35,thin] coordinates {(1,0.000) (2,-0.000) (3,0.000) (4,0.000) (5,0.000) (6,0.000) (7,-0.000) (8,0.000) (9,-0.000) (10,-0.000) (11,0.000) (12,-0.000) (13,0.000) (14,-0.000) (15,0.000) (16,0.000) (17,-0.000) (18,-0.000) (19,-0.000) (20,-0.000) (21,-0.000) (22,0.000) (23,0.000) (24,0.000) (25,0.000) (26,0.000) (27,-0.000) (28,0.000) (29,-0.000) (30,0.000) (31,0.000) (32,-0.000) (33,0.000) (34,-0.000) (35,0.000) (36,-0.000) (37,-0.000) (38,0.000) (39,-0.000) (40,-0.000) (41,-0.000) (42,-0.000) (43,0.000) (44,-0.000) (45,0.000) (46,0.000) (47,-0.000) (48,-0.000) (49,-0.000) (50,-0.000) (51,0.000) (52,-0.000) (53,0.000) (54,0.000) (55,-0.000) (56,-0.000) (57,-0.000) (58,-0.000) (59,-0.000) (60,0.000) (61,0.000) (62,-0.000) (63,-0.000) (64,0.000) (65,-0.000) (66,-0.000) (67,0.000) (68,0.000) (69,0.000) (70,-0.000) (71,-0.000) (72,-0.000) (73,0.000) (74,-0.000) (75,0.000) (76,-0.000) (77,-0.000) (78,-0.000) (79,-0.000) (80,-0.000) (81,0.000) (82,0.000) (83,-0.000) (84,-0.000) (85,0.000) (86,0.000) (87,-0.000) (88,-0.000) (89,0.000) (90,0.000) (91,0.000) (92,-0.000) (93,-0.000) (94,-0.000) (95,-0.000) (96,-0.000) (97,0.000) (98,0.000) (99,-0.000) (100,-0.000) (101,0.000) (102,0.000) (103,0.000) (104,0.000) (105,-0.000) (106,-0.000) (107,-0.000) (108,-0.000) (109,0.000) (110,-0.000) (111,-0.000) (112,-0.000) (113,-0.000) (114,-0.000) (115,-0.000) (116,-0.000) (117,0.000) (118,-0.000) (119,-0.000) (120,-0.000) (121,0.000) (122,-0.000) (123,-0.000) (124,-0.000) (125,0.000) (126,-0.000) (127,-0.000) (128,0.000) (129,-0.000) (130,-0.000) (131,0.000) (132,-0.000) (133,-0.000) (134,-0.000) (135,0.000) (136,-0.000) (137,-0.000) (138,-0.000) (139,-0.000) (140,-0.000) (141,0.000) (142,0.000) (143,0.000) (144,0.000) (145,-0.000) (146,0.000) (147,0.000) (148,0.000) (149,-0.000) (150,0.000)};
\addplot[sysblue,very thick] coordinates {(1,0.000) (3,0.000) (5,0.000) (7,0.000) (9,0.000) (11,0.000) (13,0.000) (15,0.000) (17,0.000) (19,0.000) (21,-0.000) (23,-0.000) (25,-0.000) (27,-0.000) (29,-0.000) (31,-0.000) (33,-0.000) (35,-0.000) (37,-0.000) (39,-0.000) (41,-0.000) (43,-0.000) (45,-0.000) (47,-0.000) (49,-0.000) (51,-0.000) (53,-0.000) (55,-0.000) (57,-0.000) (59,-0.000) (61,-0.000) (63,-0.000) (65,-0.000) (67,-0.000) (69,-0.000) (71,-0.000) (73,-0.000) (75,-0.000) (77,-0.000) (79,-0.000) (81,-0.000) (83,-0.000) (85,-0.000) (87,-0.000) (89,-0.000) (91,-0.000) (93,-0.000) (95,-0.000) (97,-0.000) (99,-0.000) (101,-0.000) (103,-0.000) (105,-0.000) (107,-0.000) (109,-0.000) (111,-0.000) (113,-0.000) (115,-0.000) (117,-0.000) (119,-0.000) (121,-0.000) (123,-0.000) (125,-0.000) (127,-0.000) (129,-0.000) (131,-0.000) (133,-0.000) (135,-0.000) (137,-0.000) (139,-0.000) (141,-0.000) (143,-0.000) (145,-0.000) (147,-0.000) (149,-0.000)};
\nextgroupplot[ylabel={advantage std},ymin=0,ymax=1.1]
\addplot[sysblue!35,thin] coordinates {(1,0.667) (2,0.667) (3,0.816) (4,0.943) (5,0.882) (6,0.816) (7,0.882) (8,0.943) (9,0.882) (10,0.882) (11,0.943) (12,0.943) (13,0.943) (14,0.745) (15,0.943) (16,0.943) (17,0.943) (18,0.882) (19,0.816) (20,0.882) (21,0.816) (22,0.882) (23,0.745) (24,0.943) (25,0.882) (26,0.745) (27,0.816) (28,0.882) (29,0.943) (30,0.745) (31,0.943) (32,0.577) (33,0.816) (34,0.667) (35,0.882) (36,0.667) (37,0.745) (38,0.882) (39,0.882) (40,0.882) (41,0.882) (42,0.745) (43,0.745) (44,0.816) (45,0.816) (46,0.882) (47,0.816) (48,0.882) (49,0.816) (50,0.816) (51,0.816) (52,0.943) (53,0.882) (54,0.745) (55,0.943) (56,0.816) (57,0.666) (58,0.816) (59,0.882) (60,0.745) (61,0.577) (62,0.816) (63,0.882) (64,0.745) (65,0.943) (66,0.745) (67,0.816) (68,0.816) (69,0.816) (70,0.745) (71,0.667) (72,0.745) (73,0.577) (74,0.816) (75,0.577) (76,0.816) (77,0.667) (78,0.882) (79,0.816) (80,0.667) (81,0.745) (82,0.745) (83,0.667) (84,0.882) (85,0.745) (86,0.667) (87,0.816) (88,0.667) (89,0.816) (90,0.577) (91,0.816) (92,0.816) (93,0.333) (94,0.577) (95,0.666) (96,0.816) (97,0.577) (98,0.745) (99,0.667) (100,0.667) (101,0.816) (102,0.745) (103,0.667) (104,0.816) (105,0.745) (106,0.816) (107,0.882) (108,0.816) (109,0.471) (110,0.667) (111,0.816) (112,0.745) (113,0.667) (114,0.745) (115,0.816) (116,0.577) (117,0.667) (118,0.577) (119,0.577) (120,0.816) (121,0.471) (122,0.577) (123,0.333) (124,0.666) (125,0.882) (126,0.667) (127,0.667) (128,0.745) (129,0.667) (130,0.577) (131,0.577) (132,0.666) (133,0.667) (134,0.577) (135,0.333) (136,0.667) (137,0.745) (138,0.666) (139,0.577) (140,0.666) (141,0.816) (142,0.667) (143,0.577) (144,0.577) (145,0.471) (146,0.577) (147,0.667) (148,0.745) (149,0.577) (150,0.667)};
\addplot[sysblue,very thick] coordinates {(1,0.667) (3,0.716) (5,0.795) (7,0.810) (9,0.833) (11,0.847) (13,0.862) (15,0.860) (17,0.869) (19,0.867) (21,0.875) (23,0.883) (25,0.883) (27,0.876) (29,0.876) (31,0.869) (33,0.844) (35,0.837) (37,0.814) (39,0.817) (41,0.820) (43,0.813) (45,0.804) (47,0.811) (49,0.804) (51,0.802) (53,0.823) (55,0.830) (57,0.834) (59,0.830) (61,0.808) (63,0.819) (65,0.821) (67,0.815) (69,0.811) (71,0.800) (73,0.775) (75,0.761) (77,0.761) (79,0.761) (81,0.765) (83,0.751) (85,0.748) (87,0.744) (89,0.736) (91,0.735) (93,0.727) (95,0.719) (97,0.715) (99,0.700) (101,0.704) (103,0.704) (105,0.701) (107,0.711) (109,0.702) (111,0.706) (113,0.719) (115,0.735) (117,0.728) (119,0.715) (121,0.705) (123,0.680) (125,0.679) (127,0.661) (129,0.667) (131,0.651) (133,0.647) (135,0.614) (137,0.623) (139,0.627) (141,0.637) (143,0.654) (145,0.629) (147,0.624) (149,0.620)};
\nextgroupplot[ylabel={entropy (nats)},ymin=0,ymax=0.7,xlabel={GRPO step}]
\addplot[sysblue!35,thin] coordinates {(1,0.527) (2,0.499) (3,0.478) (4,0.480) (5,0.436) (6,0.425) (7,0.422) (8,0.421) (9,0.353) (10,0.342) (11,0.412) (12,0.357) (13,0.439) (14,0.303) (15,0.372) (16,0.308) (17,0.347) (18,0.346) (19,0.369) (20,0.320) (21,0.312) (22,0.303) (23,0.334) (24,0.317) (25,0.286) (26,0.296) (27,0.297) (28,0.288) (29,0.275) (30,0.304) (31,0.280) (32,0.197) (33,0.224) (34,0.203) (35,0.293) (36,0.214) (37,0.194) (38,0.263) (39,0.276) (40,0.298) (41,0.224) (42,0.237) (43,0.235) (44,0.280) (45,0.224) (46,0.274) (47,0.298) (48,0.278) (49,0.184) (50,0.196) (51,0.197) (52,0.239) (53,0.172) (54,0.249) (55,0.261) (56,0.237) (57,0.163) (58,0.223) (59,0.250) (60,0.175) (61,0.169) (62,0.186) (63,0.226) (64,0.216) (65,0.229) (66,0.155) (67,0.126) (68,0.226) (69,0.157) (70,0.156) (71,0.198) (72,0.173) (73,0.119) (74,0.214) (75,0.086) (76,0.224) (77,0.146) (78,0.194) (79,0.159) (80,0.150) (81,0.171) (82,0.160) (83,0.149) (84,0.132) (85,0.206) (86,0.151) (87,0.102) (88,0.167) (89,0.161) (90,0.118) (91,0.157) (92,0.089) (93,0.080) (94,0.135) (95,0.145) (96,0.179) (97,0.128) (98,0.161) (99,0.139) (100,0.094) (101,0.144) (102,0.115) (103,0.132) (104,0.157) (105,0.125) (106,0.111) (107,0.110) (108,0.136) (109,0.104) (110,0.159) (111,0.174) (112,0.165) (113,0.129) (114,0.102) (115,0.157) (116,0.086) (117,0.127) (118,0.091) (119,0.100) (120,0.125) (121,0.080) (122,0.073) (123,0.064) (124,0.111) (125,0.169) (126,0.105) (127,0.099) (128,0.093) (129,0.116) (130,0.082) (131,0.087) (132,0.094) (133,0.144) (134,0.105) (135,0.077) (136,0.123) (137,0.079) (138,0.085) (139,0.104) (140,0.104) (141,0.132) (142,0.084) (143,0.070) (144,0.035) (145,0.064) (146,0.065) (147,0.107) (148,0.085) (149,0.078) (150,0.091)};
\addplot[sysblue,very thick] coordinates {(1,0.527) (3,0.501) (5,0.484) (7,0.467) (9,0.449) (11,0.436) (13,0.430) (15,0.418) (17,0.407) (19,0.402) (21,0.387) (23,0.370) (25,0.354) (27,0.342) (29,0.331) (31,0.323) (33,0.304) (35,0.295) (37,0.283) (39,0.274) (41,0.268) (43,0.260) (45,0.255) (47,0.254) (49,0.249) (51,0.239) (53,0.239) (55,0.240) (57,0.239) (59,0.236) (61,0.227) (63,0.224) (65,0.221) (67,0.206) (69,0.203) (71,0.201) (73,0.195) (75,0.184) (77,0.183) (79,0.177) (81,0.175) (83,0.170) (85,0.165) (87,0.164) (89,0.161) (91,0.157) (93,0.151) (95,0.150) (97,0.147) (99,0.144) (101,0.140) (103,0.137) (105,0.134) (107,0.132) (109,0.128) (111,0.131) (113,0.137) (115,0.136) (117,0.131) (119,0.126) (121,0.124) (123,0.119) (125,0.119) (127,0.118) (129,0.116) (131,0.108) (133,0.105) (135,0.101) (137,0.101) (139,0.101) (141,0.102) (143,0.103) (145,0.094) (147,0.093) (149,0.090)};
\nextgroupplot[ylabel={zero-variance fraction},ymin=0,ymax=1.05,xlabel={GRPO step}]
\addplot[sysblue!35,thin] coordinates {(1,0.500) (2,0.500) (3,0.250) (4,0.000) (5,0.125) (6,0.250) (7,0.125) (8,0.000) (9,0.125) (10,0.125) (11,0.000) (12,0.000) (13,0.000) (14,0.375) (15,0.000) (16,0.000) (17,0.000) (18,0.125) (19,0.250) (20,0.125) (21,0.250) (22,0.125) (23,0.375) (24,0.000) (25,0.125) (26,0.375) (27,0.250) (28,0.125) (29,0.000) (30,0.375) (31,0.000) (32,0.625) (33,0.250) (34,0.500) (35,0.125) (36,0.500) (37,0.375) (38,0.125) (39,0.125) (40,0.125) (41,0.125) (42,0.375) (43,0.375) (44,0.250) (45,0.250) (46,0.125) (47,0.250) (48,0.125) (49,0.250) (50,0.250) (51,0.250) (52,0.000) (53,0.125) (54,0.375) (55,0.000) (56,0.250) (57,0.500) (58,0.250) (59,0.125) (60,0.375) (61,0.625) (62,0.250) (63,0.125) (64,0.375) (65,0.000) (66,0.375) (67,0.250) (68,0.250) (69,0.250) (70,0.375) (71,0.500) (72,0.375) (73,0.625) (74,0.250) (75,0.625) (76,0.250) (77,0.500) (78,0.125) (79,0.250) (80,0.500) (81,0.375) (82,0.375) (83,0.500) (84,0.125) (85,0.375) (86,0.500) (87,0.250) (88,0.500) (89,0.250) (90,0.625) (91,0.250) (92,0.250) (93,0.875) (94,0.625) (95,0.500) (96,0.250) (97,0.625) (98,0.375) (99,0.500) (100,0.500) (101,0.250) (102,0.375) (103,0.500) (104,0.250) (105,0.375) (106,0.250) (107,0.125) (108,0.250) (109,0.750) (110,0.500) (111,0.250) (112,0.375) (113,0.500) (114,0.375) (115,0.250) (116,0.625) (117,0.500) (118,0.625) (119,0.625) (120,0.250) (121,0.750) (122,0.625) (123,0.875) (124,0.500) (125,0.125) (126,0.500) (127,0.500) (128,0.375) (129,0.500) (130,0.625) (131,0.625) (132,0.500) (133,0.500) (134,0.625) (135,0.875) (136,0.500) (137,0.375) (138,0.500) (139,0.625) (140,0.500) (141,0.250) (142,0.500) (143,0.625) (144,0.625) (145,0.750) (146,0.625) (147,0.500) (148,0.375) (149,0.625) (150,0.500)};
\addplot[sysblue,very thick] coordinates {(1,0.500) (3,0.417) (5,0.275) (7,0.250) (9,0.208) (11,0.182) (13,0.154) (15,0.158) (17,0.140) (19,0.145) (21,0.131) (23,0.119) (25,0.119) (27,0.131) (29,0.131) (31,0.144) (33,0.188) (35,0.200) (37,0.244) (39,0.237) (41,0.231) (43,0.244) (45,0.263) (47,0.250) (49,0.263) (51,0.269) (53,0.231) (55,0.219) (57,0.212) (59,0.219) (61,0.256) (63,0.237) (65,0.231) (67,0.244) (69,0.250) (71,0.269) (73,0.312) (75,0.338) (77,0.338) (79,0.338) (81,0.331) (83,0.356) (85,0.362) (87,0.369) (89,0.381) (91,0.381) (93,0.388) (95,0.400) (97,0.406) (99,0.431) (101,0.425) (103,0.425) (105,0.431) (107,0.412) (109,0.425) (111,0.419) (113,0.406) (115,0.381) (117,0.394) (119,0.412) (121,0.425) (123,0.456) (125,0.456) (127,0.487) (129,0.481) (131,0.506) (133,0.512) (135,0.556) (137,0.544) (139,0.537) (141,0.525) (143,0.506) (145,0.544) (147,0.550) (149,0.556)};
\end{groupplot}
\end{tikzpicture}
\caption{Advanced metrics of the baseline (Chapter~37) GRPO run on the toy task, seed 0,
measured on CPU. The advantage mean is zero by construction (a free correctness check). The
advantage std and the zero-variance fraction mirror each other: as more prompts are solved
by all 8 rollouts, fewer groups carry signal. Entropy falls as the policy sharpens.}
\label{fig:grpo2:baseline}
\end{figure}

% ======================================================================
\section{Stabilising updates with clipped policy ratios}\label{sec:grpo2:clip}

\subsection{Where the ratio comes from}\label{sec:grpo2:ratio}

The policy-gradient estimator of \cref{eq:grpo:pg} assumes the rollouts were sampled from the
\emph{current} policy $\pitheta$. After one optimiser step that is no longer true: the
rollouts came from the snapshot $\piold$ that generated them. If we want to reuse a batch
of expensive rollouts for several updates, or if the sampler runs a few steps behind the
trainer (\cref{ch:rlsys}), we must correct for the mismatch. Importance sampling does this:
\begin{equation}
\E_{y\sim\pitheta}\!\big[A\,\nabla\log\pitheta(y)\big]
=\E_{y\sim\piold}\!\Big[\frac{\pitheta(y)}{\piold(y)}A\,\nabla\log\pitheta(y)\Big]
=\E_{y\sim\piold}\!\big[A\,\nabla\rho(y)\big],\qquad \rho=\frac{\pitheta(y)}{\piold(y)},
\end{equation}
so maximising the \emph{surrogate} $\E_{\piold}[\rho A]$ has the right gradient. The ratio
equals 1 on the first update after sampling (which is why the summed-log-prob loss of
\cref{ch:grpo} gives the same gradient) and drifts away from 1 with every further update on
the same data. Unchecked, the surrogate can be pushed arbitrarily high by making $\rho$ huge
on positive-advantage samples. The estimate is then dominated by a few samples and the
policy leaves the region where the old samples say anything about it.

PPO's answer is to clip the ratio to a trust region $[1-\varepsilon,1+\varepsilon]$ and take
the pessimistic minimum:
\begin{equation}
\ell^{\text{clip}}(\rho,A)=\min\big(\rho A,\ \clip(\rho,1-\varepsilon,1+\varepsilon)A\big).
\label{eq:grpo2:clip}
\end{equation}
\Cref{fig:grpo2:clipobj} shows what the minimum does. For a good sample ($A>0$) the
objective stops increasing once $\rho>1+\varepsilon$: there is no reward for making it
\emph{even} more likely in this update, so its gradient is zero. For a bad sample ($A<0$) the
objective stops improving once $\rho<1-\varepsilon$. In the opposite directions, a good sample
whose probability has \emph{dropped} or a bad sample whose probability has \emph{risen}, the
unclipped term is the smaller one and the gradient stays on to correct the mistake.

\begin{figure}[ht]
\centering
\begin{tikzpicture}
\begin{groupplot}[group style={group size=2 by 1,horizontal sep=1.6cm},width=0.47\linewidth,height=4.4cm,
  xmin=0,xmax=2,xlabel={ratio $\rho=\pitheta/\piold$},grid=major,grid style={gray!15},
  tick label style={font=\scriptsize},label style={font=\small},title style={font=\small}]
\nextgroupplot[title={$A=+1$},ymin=-0.1,ymax=2.1,ylabel={$\ell^{\text{clip}}$}]
\addplot[gray,dashed,domain=0:2] {x};
\addplot[very thick,keygreen,domain=0:1.2] {x};
\addplot[very thick,keygreen,domain=1.2:2] {1.2};
\addplot[fill=keygreen!12,draw=none] coordinates {(1.2,-0.1)(2,-0.1)(2,2.1)(1.2,2.1)} \closedcycle;
\node[font=\scriptsize,align=center,keygreen!70!black] at (axis cs:1.6,0.55) {flat:\\no gradient};
\node[font=\scriptsize,gray] at (axis cs:1.75,1.9) {$\rho A$};
\nextgroupplot[title={$A=-1$},ymin=-2.1,ymax=0.1]
\addplot[gray,dashed,domain=0:2] {-x};
\addplot[very thick,warnred,domain=0:0.8] {-0.8};
\addplot[very thick,warnred,domain=0.8:2] {-x};
\addplot[fill=warnred!10,draw=none] coordinates {(0,-2.1)(0.8,-2.1)(0.8,0.1)(0,0.1)} \closedcycle;
\node[font=\scriptsize,align=center,warnred!70!black] at (axis cs:0.4,-1.4) {flat:\\no gradient};
\end{groupplot}
\end{tikzpicture}
\caption{The clipped surrogate of \cref{eq:grpo2:clip} with $\varepsilon=0.2$ (solid) versus
the unclipped $\rho A$ (dashed). Shaded regions are where the clip binds and the sample's
gradient is zero: moving further in the direction the advantage wants brings no further
credit within this update. Moving the wrong way is never clipped.}
\label{fig:grpo2:clipobj}
\end{figure}

A worked example on the four rollouts of \cref{fig:grpo:step} makes this concrete. Suppose
that after some updates the current policy's sequence log-probabilities are
$(-2.269,-3.059,-1.299,-3.400)$, while the snapshot that sampled them had
$(-2.469,-2.659,-1.299,-3.100)$. Then

\begin{center}\small
\begin{tabular}{@{}lrrrr@{}}
\toprule
rollout & $A_i$ & $\rho_i=e^{\Delta\log\pi}$ & $\clip(\rho_i,0.8,1.2)$ & $\min(\rho A,\clip(\rho)A)$\\
\midrule
\verb|\boxed{7}|              & $+0.866$ & 1.221 & 1.200 & $1.039$ (clipped)\\
\verb|<think>..\boxed{7}| & $+0.866$ & 0.670 & 0.800 & $0.581$ (unclipped)\\
\verb|7|                      & $-0.866$ & 1.000 & 1.000 & $-0.866$\\
\verb|\boxed{8}|              & $-0.866$ & 0.741 & 0.800 & $-0.693$ (clipped)\\
\bottomrule
\end{tabular}
\end{center}

The clipped loss is $-\frac14\sum\ell^{\text{clip}}_i=-0.0152$, against $-0.0327$
unclipped. Two rollouts no longer contribute gradient. The first boxed answer has already
been made 22\% more likely, and the wrong boxed answer 26\% less likely, so both are
``done'' for this batch. The second correct answer became \emph{less} likely, an error
the unclipped branch keeps correcting.

\begin{lstlisting}[caption={Policy ratios and the clipped loss (sequence level, as in this section).},label={lst:grpo2:clip}]
def clipped_policy_loss(new_logps, old_logps, advantages,
                        eps_low=0.2, eps_high=0.2):
    ratio = torch.exp(new_logps - old_logps.detach())
    adv = advantages.detach()
    unclipped = ratio * adv
    clipped = torch.clamp(ratio, 1 - eps_low, 1 + eps_high) * adv
    obj = torch.minimum(unclipped, clipped)
    clip_frac = (clipped < unclipped).float().mean().item()  # clip binds
    return -obj.mean(), ratio.mean().item(), clip_frac
\end{lstlisting}

\paragraph{Choosing $\varepsilon$.} Classic PPO for RLHF uses $\varepsilon\in[0.1,0.2]$,
which is aggressive clipping and small per-update changes. Reasoning-RL recipes span a wide
range, and some large runs reportedly used very loose clipping, where $\varepsilon$ is so large
that only extreme ratios are touched. The right value depends on how many times each rollout
is reused and how stale the sampler is. Two refinements are common. \emph{Asymmetric} clipping
($\varepsilon_{\text{low}}<\varepsilon_{\text{high}}$, DAPO's ``clip-higher'') lets rare good
tokens grow faster than common ones shrink, which counteracts entropy collapse.
\emph{Sequence-level} versus \emph{token-level} ratios is the other choice: this section's
listing clips the ratio of whole-sequence probabilities, while the GRPO paper clips every
token's ratio $\rho_{i,t}$ separately (\cref{eq:grpo:full}). Token-level ratios are numerically
tamer, because a sequence ratio is a product of thousands of per-token ratios. GSPO
(\cref{sec:grpo2:more}) goes back to the sequence level but takes a geometric mean.

\begin{keyidea}[log the ratio, not just the loss]
Two cheap diagnostics come with clipping: the mean ratio (should hover near 1; persistent
drift means the policy is moving faster than the batch can support) and the
\emph{clip fraction} (fraction of tokens where the clip binds). A clip fraction near 0
means clipping is idle and $\varepsilon$ or the reuse can be increased. One above 0.2--0.3
means most of the batch is being ignored and updates are too aggressive.
\end{keyidea}

\subsection{Training with clipped ratios: three snapshots of one policy}\label{sec:grpo2:clipstrain}

Clipping only does something if the same rollouts are seen by more than one optimiser step,
so implementing it forces an explicit choice about \emph{which version} of the policy each
quantity comes from. A GRPO implementation juggles three snapshots (\cref{fig:grpo2:snapshots}):
\begin{itemize}
\item $\pitheta$, the live weights, changed by every optimiser step;
\item $\piold$, the weights that generated the current rollout batch, refreshed every time a
new batch is sampled. Its log-probabilities are computed once per batch under \code{no\_grad}
and stored as numbers, not as a model copy;
\item $\piref$, the frozen anchor for the KL term (\cref{sec:grpo2:kl}), usually the
initial model, refreshed rarely if ever.
\end{itemize}

\begin{figure}[ht]
\centering
\resizebox{\linewidth}{!}{%
\begin{tikzpicture}[>={Stealth[length=2mm]},font=\small]
\draw[->,thick] (0,0)--(13.2,0) node[right,font=\scriptsize]{optimiser steps};
\foreach \x in {0.5,1.2,...,12.5} \draw (\x,-0.07)--(\x,0.07);
\node[anchor=east,font=\scriptsize] at (-0.1,0) {$\pitheta$};
% batches
\foreach \b/\xa/\xb in {1/0.3/4.4,2/4.5/8.6,3/8.7/12.8} {
  \draw[fill=fcompute!15,draw=fcompute,rounded corners=2pt] (\xa,0.45) rectangle (\xb,1.05);
  \node[font=\scriptsize] at ({(\xa+\xb)/2},0.75) {batch \b: $\piold$ fixed, $k$ updates};
  \draw[->,fcompute] (\xa+0.05,1.05) -- (\xa+0.05,1.45);
}
\node[anchor=east,font=\scriptsize] at (-0.1,0.75) {$\piold$};
\node[font=\scriptsize,fcompute] at (0.6,1.65) {sample};
\node[font=\scriptsize,fcompute] at (4.8,1.65) {sample};
\node[font=\scriptsize,fcompute] at (9.0,1.65) {sample};
\draw[fill=gray!15,draw=gray,rounded corners=2pt] (0.3,-0.55) rectangle (12.8,-1.1);
\node[font=\scriptsize] at (6.5,-0.82) {$\piref$: frozen initial model (or refreshed every few hundred steps)};
\node[anchor=east,font=\scriptsize] at (-0.1,-0.82) {$\piref$};
\end{tikzpicture}}
\caption{Three snapshots of one policy. Each rollout batch fixes $\piold$ (stored as
log-probabilities). The live $\pitheta$ takes several minibatch steps on that batch, and the
ratio $\pitheta/\piold$ measures how far it has moved; clipping bounds that. $\piref$
anchors the whole run through the KL term.}
\label{fig:grpo2:snapshots}
\end{figure}

DeepSeek-R1-style training samples one large rollout pool per iteration (thousands of
completions), then walks through it in many minibatches, each an optimiser step, before
sampling again. Later minibatches therefore see ratios well away from 1, and clipping does
real work. Our trainer exposes the same structure through two knobs of
\code{GRPOConfig}: \code{minibatches} (how many optimiser steps one batch is split into) and
\code{inner\_epochs} (how many passes over the batch). With \code{inner\_epochs=1} and
\code{minibatches=1}, every ratio is exactly 1, the clipped loss reduces to
\cref{ch:grpo}'s gradient, and the clip never binds. (In that setting the logged loss is also
exactly zero, a computational artefact, because $\rho\equiv1$ and the advantages are
mean-zero.)

\begin{systemslens}[rollout reuse is a throughput knob]
Reusing each rollout batch for $k$ optimiser steps divides the sampling cost per step by
$k$. Since sampling dominates (\cref{sec:grpo:loop}), $k=2$--4 can nearly halve wall-clock time
per update. The price is off-policyness, and clipping is what makes it affordable. The same
mathematics covers \emph{asynchronous} RL, where the sampler generates batch $n{+}1$ with
weights from step $n-s$ while the trainer is at step $n$. To the loss, a stale sampler looks
exactly like rollout reuse, and the ratio and clip absorb it, up to a point (\cref{ch:rlsys}).
\end{systemslens}

\paragraph{What the toy shows.} We ran two reuse settings from the same SFT start
(\cref{fig:grpo2:clipruns}). With \emph{moderate} reuse, four full passes over each rollout
batch at the base learning rate, the unclipped run stalled at a reward near 0.5 and ended
with greedy accuracy 0.62 and 0.50 (two seeds). The clipped run ($\varepsilon=0.2$) matched
the no-reuse baseline: 0.88 and 0.95, against 0.89 and 0.96. Its clip fraction was only
2--3\%. A few runaway tokens per batch are enough to derail the unclipped run, and these
are precisely the ones clipping silences. With \emph{aggressive} reuse, eight optimiser steps
per batch at three times the learning rate, \emph{both} runs collapsed: entropy fell to
0.01--0.06, every group became uniformly wrong (zero-variance fraction $\approx1$), and greedy
accuracy ended at 0.07--0.16. The lesson is about what clipping does and does not bound. It
bounds the drift \emph{within} one batch relative to $\piold$. Every new batch opens a new
trust region, so clipping is no substitute for a sensible learning rate.

\begin{figure}[ht]
\centering
\begin{tikzpicture}
\begin{groupplot}[group style={group size=2 by 1,horizontal sep=1.5cm},width=0.48\linewidth,height=4.3cm,
  xmin=0,xmax=150,xlabel={GRPO step},grid=major,grid style={gray!15},tick label style={font=\scriptsize},label style={font=\small}]
\nextgroupplot[ylabel={reward (train)},ymin=0,ymax=1.0,legend to name=cliplegend,legend columns=3,legend style={font=\scriptsize,draw=none}]
\addplot[gray,very thick] coordinates {(1,0.102) (4,0.195) (7,0.241) (10,0.293) (13,0.332) (16,0.391) (19,0.437) (22,0.483) (25,0.504) (28,0.522) (31,0.528) (34,0.545) (37,0.574) (40,0.590) (43,0.598) (46,0.609) (49,0.630) (52,0.624) (55,0.641) (58,0.643) (61,0.653) (64,0.654) (67,0.652) (70,0.649) (73,0.655) (76,0.654) (79,0.658) (82,0.689) (85,0.701) (88,0.737) (91,0.756) (94,0.781) (97,0.772) (100,0.781) (103,0.782) (106,0.786) (109,0.796) (112,0.804) (115,0.822) (118,0.835) (121,0.834) (124,0.824) (127,0.824) (130,0.809) (133,0.799) (136,0.806) (139,0.817) (142,0.814) (145,0.837) (148,0.836) (150,0.831)};
\addlegendentry{baseline (no reuse)}
\addplot[warnred,very thick] coordinates {(1,0.102) (4,0.328) (7,0.348) (10,0.392) (13,0.377) (16,0.435) (19,0.457) (22,0.466) (25,0.441) (28,0.486) (31,0.457) (34,0.456) (37,0.459) (40,0.500) (43,0.480) (46,0.518) (49,0.503) (52,0.503) (55,0.496) (58,0.487) (61,0.473) (64,0.504) (67,0.513) (70,0.512) (73,0.523) (76,0.499) (79,0.473) (82,0.466) (85,0.457) (88,0.450) (91,0.508) (94,0.550) (97,0.552) (100,0.535) (103,0.555) (106,0.506) (109,0.463) (112,0.486) (115,0.506) (118,0.529) (121,0.564) (124,0.591) (127,0.588) (130,0.594) (133,0.587) (136,0.554) (139,0.527) (142,0.494) (145,0.530) (148,0.523) (150,0.526)};
\addlegendentry{4 passes, no clip}
\addplot[keygreen,very thick] coordinates {(1,0.102) (4,0.242) (7,0.335) (10,0.398) (13,0.427) (16,0.512) (19,0.578) (22,0.605) (25,0.631) (28,0.659) (31,0.624) (34,0.632) (37,0.656) (40,0.685) (43,0.698) (46,0.725) (49,0.750) (52,0.746) (55,0.740) (58,0.737) (61,0.753) (64,0.761) (67,0.771) (70,0.760) (73,0.767) (76,0.753) (79,0.749) (82,0.760) (85,0.778) (88,0.794) (91,0.820) (94,0.828) (97,0.802) (100,0.784) (103,0.791) (106,0.791) (109,0.783) (112,0.822) (115,0.830) (118,0.818) (121,0.805) (124,0.816) (127,0.807) (130,0.805) (133,0.809) (136,0.786) (139,0.766) (142,0.770) (145,0.790) (148,0.792) (150,0.811)};
\addlegendentry{4 passes, clip $\varepsilon=0.2$}
\addplot[warnred!50,dashed,very thick] coordinates {(1,0.102) (4,0.281) (7,0.251) (10,0.255) (13,0.218) (16,0.208) (19,0.161) (22,0.160) (25,0.131) (28,0.129) (31,0.131) (34,0.118) (37,0.111) (40,0.118) (43,0.130) (46,0.140) (49,0.169) (52,0.144) (55,0.130) (58,0.111) (61,0.086) (64,0.074) (67,0.083) (70,0.078) (73,0.086) (76,0.095) (79,0.088) (82,0.099) (85,0.099) (88,0.094) (91,0.090) (94,0.133) (97,0.117) (100,0.133) (103,0.133) (106,0.150) (109,0.108) (112,0.100) (115,0.083) (118,0.096) (121,0.075) (124,0.062) (127,0.071) (130,0.067) (133,0.056) (136,0.068) (139,0.084) (142,0.080) (145,0.101) (148,0.098) (150,0.111)};
\addlegendentry{8 steps/batch, lr$\times3$, no clip}
\addplot[keygreen!60!black,dashed,very thick] coordinates {(1,0.102) (4,0.293) (7,0.271) (10,0.278) (13,0.282) (16,0.281) (19,0.250) (22,0.245) (25,0.236) (28,0.222) (31,0.229) (34,0.226) (37,0.233) (40,0.247) (43,0.246) (46,0.225) (49,0.235) (52,0.211) (55,0.185) (58,0.166) (61,0.177) (64,0.167) (67,0.201) (70,0.205) (73,0.227) (76,0.221) (79,0.221) (82,0.213) (85,0.214) (88,0.204) (91,0.205) (94,0.196) (97,0.166) (100,0.156) (103,0.148) (106,0.166) (109,0.191) (112,0.209) (115,0.207) (118,0.211) (121,0.199) (124,0.179) (127,0.146) (130,0.154) (133,0.138) (136,0.121) (139,0.121) (142,0.146) (145,0.130) (148,0.152) (150,0.161)};
\addlegendentry{8 steps/batch, lr$\times3$, clip}
\nextgroupplot[ylabel={entropy (nats)},ymin=0,ymax=0.6]
\addplot[gray,very thick] coordinates {(1,0.532) (4,0.500) (7,0.467) (10,0.446) (13,0.435) (16,0.409) (19,0.384) (22,0.365) (25,0.351) (28,0.331) (31,0.319) (34,0.292) (37,0.275) (40,0.263) (43,0.250) (46,0.243) (49,0.246) (52,0.240) (55,0.230) (58,0.228) (61,0.219) (64,0.210) (67,0.203) (70,0.194) (73,0.185) (76,0.178) (79,0.173) (82,0.171) (85,0.171) (88,0.163) (91,0.160) (94,0.153) (97,0.151) (100,0.146) (103,0.144) (106,0.140) (109,0.137) (112,0.137) (115,0.132) (118,0.127) (121,0.120) (124,0.120) (127,0.114) (130,0.110) (133,0.113) (136,0.113) (139,0.110) (142,0.109) (145,0.102) (148,0.096) (150,0.096)};
\addplot[warnred,very thick] coordinates {(1,0.532) (4,0.438) (7,0.384) (10,0.343) (13,0.318) (16,0.278) (19,0.228) (22,0.191) (25,0.172) (28,0.154) (31,0.144) (34,0.126) (37,0.123) (40,0.114) (43,0.104) (46,0.093) (49,0.093) (52,0.088) (55,0.082) (58,0.079) (61,0.079) (64,0.077) (67,0.072) (70,0.068) (73,0.068) (76,0.068) (79,0.073) (82,0.073) (85,0.073) (88,0.073) (91,0.067) (94,0.060) (97,0.064) (100,0.065) (103,0.067) (106,0.068) (109,0.068) (112,0.065) (115,0.061) (118,0.055) (121,0.053) (124,0.058) (127,0.052) (130,0.055) (133,0.054) (136,0.052) (139,0.049) (142,0.057) (145,0.052) (148,0.056) (150,0.059)};
\addplot[keygreen,very thick] coordinates {(1,0.532) (4,0.457) (7,0.402) (10,0.365) (13,0.334) (16,0.285) (19,0.235) (22,0.203) (25,0.179) (28,0.160) (31,0.163) (34,0.151) (37,0.140) (40,0.132) (43,0.126) (46,0.116) (49,0.110) (52,0.105) (55,0.104) (58,0.103) (61,0.097) (64,0.092) (67,0.091) (70,0.087) (73,0.083) (76,0.085) (79,0.086) (82,0.085) (85,0.079) (88,0.073) (91,0.064) (94,0.057) (97,0.055) (100,0.053) (103,0.056) (106,0.056) (109,0.056) (112,0.051) (115,0.053) (118,0.052) (121,0.051) (124,0.051) (127,0.054) (130,0.051) (133,0.050) (136,0.056) (139,0.059) (142,0.059) (145,0.056) (148,0.055) (150,0.051)};
\addplot[warnred!50,dashed,very thick] coordinates {(1,0.532) (4,0.305) (7,0.218) (10,0.179) (13,0.160) (16,0.123) (19,0.099) (22,0.096) (25,0.099) (28,0.090) (31,0.080) (34,0.064) (37,0.053) (40,0.043) (43,0.039) (46,0.038) (49,0.046) (52,0.048) (55,0.043) (58,0.048) (61,0.049) (64,0.044) (67,0.047) (70,0.046) (73,0.042) (76,0.037) (79,0.034) (82,0.026) (85,0.023) (88,0.020) (91,0.019) (94,0.017) (97,0.016) (100,0.019) (103,0.018) (106,0.018) (109,0.018) (112,0.019) (115,0.019) (118,0.018) (121,0.019) (124,0.021) (127,0.019) (130,0.019) (133,0.024) (136,0.025) (139,0.026) (142,0.029) (145,0.029) (148,0.026) (150,0.025)};
\addplot[keygreen!60!black,dashed,very thick] coordinates {(1,0.532) (4,0.355) (7,0.273) (10,0.251) (13,0.222) (16,0.174) (19,0.141) (22,0.125) (25,0.101) (28,0.091) (31,0.089) (34,0.078) (37,0.076) (40,0.082) (43,0.082) (46,0.079) (49,0.070) (52,0.061) (55,0.051) (58,0.053) (61,0.053) (64,0.059) (67,0.059) (70,0.054) (73,0.045) (76,0.043) (79,0.040) (82,0.047) (85,0.050) (88,0.055) (91,0.057) (94,0.064) (97,0.065) (100,0.065) (103,0.062) (106,0.058) (109,0.047) (112,0.046) (115,0.048) (118,0.053) (121,0.056) (124,0.061) (127,0.066) (130,0.063) (133,0.057) (136,0.055) (139,0.053) (142,0.044) (145,0.040) (148,0.038) (150,0.034)};
\end{groupplot}
\node at ($(group c1r1.south)!0.5!(group c2r1.south)+(0,-1.25)$) {\pgfplotslegendfromname{cliplegend}};
\end{tikzpicture}
\caption{Rollout reuse with and without clipping on the toy task (mean of two seeds, 15-step
moving average, measured on CPU). With four passes over each batch, the unclipped run stalls near a reward of 0.5
while the clipped run tracks the baseline. With eight
optimiser steps per batch at three times the learning rate, both runs collapse: entropy
falls towards zero and every group becomes uniformly wrong. Clipping bounds the drift
\emph{within} a batch, not the step size across batches.}
\label{fig:grpo2:clipruns}
\end{figure}

% ======================================================================
\section{Controlling drift with a KL penalty}\label{sec:grpo2:kl}

If you know classical RL you will recognise what we have so far. The update of \cref{ch:grpo}
is REINFORCE with a group baseline, and with the ratio and clip it becomes PPO without a
critic. The last ingredient of the original GRPO objective is a penalty on the
Kullback--Leibler divergence between the current policy and a fixed reference policy, usually
the model we started from:
\begin{equation}
\KL(\pitheta\,\|\,\piref)=\E_{y\sim\pitheta}\Big[\log\frac{\pitheta(y\mid x)}{\piref(y\mid x)}\Big]\ \ge 0.
\end{equation}
Clipping and KL are complementary leashes. Clipping limits how far the policy moves
\emph{within one batch}, relative to the snapshot that sampled it. The KL term limits how far it
moves \emph{over the whole run}, relative to where it started. The motivation in RLHF was
protection against reward hacking, keeping the policy in the region where the reward model is
trustworthy. With verifiable rewards there is no reward model to fool, which is one reason
recent work questions the term (below).

\subsection{Implementing the KL term}\label{sec:grpo2:klimpl}

We can only evaluate the KL on the samples we have, so we need a Monte Carlo
\emph{estimator}. Writing $r=\piref(y_t)/\pitheta(y_t)$ for a sampled token, three estimators
are in common use (\cref{fig:grpo2:klest}):
\begin{equation}
k_1=-\log r,\qquad k_2=\tfrac12(\log r)^2,\qquad k_3=r-1-\log r .
\end{equation}
$k_1$ is unbiased (its expectation under $\pitheta$ is exactly the KL) but has high variance
and is negative for half the samples. $k_2$ is always non-negative and low-variance but
biased. $k_3$, popularised by John Schulman's note on KL approximation and adopted by GRPO,
is unbiased (since $\E_{\pitheta}[r]=1$) \emph{and} always non-negative, because
$r-1\ge\log r$. GRPO applies $k_3$ per token, averages it like the policy term, and adds
$\beta\,\mathbb{D}_{i,t}$ to the per-token loss with $\beta$ around 0.001--0.05.

\begin{figure}[ht]
\centering
\begin{tikzpicture}
\begin{axis}[width=0.7\linewidth,height=4.8cm,xmin=0.2,xmax=3,ymin=-1.2,ymax=1.7,
  xlabel={$r=\piref(y_t)/\pitheta(y_t)$},ylabel={per-sample estimate},grid=major,grid style={gray!15},
  legend style={font=\scriptsize,at={(0.5,0.98)},anchor=north,legend columns=3},tick label style={font=\scriptsize}]
\addplot[very thick,warnred,domain=0.2:3,samples=80] {-ln(x)}; \addlegendentry{$k_1=-\log r$}
\addplot[very thick,napkin,domain=0.2:3,samples=80] {0.5*ln(x)^2}; \addlegendentry{$k_2$}
\addplot[very thick,sysblue,domain=0.2:3,samples=80] {x-1-ln(x)}; \addlegendentry{$k_3$}
\draw[gray,dashed] (axis cs:1,-1.2)--(axis cs:1,1.7);
\end{axis}
\end{tikzpicture}
\caption{Three per-sample KL estimators as functions of the likelihood ratio. All are zero
when policy and reference agree ($r=1$). $k_1$ goes negative whenever the policy is more
confident than the reference; $k_3$ is non-negative and has zero slope at $r=1$, so small
deviations are penalised gently.}
\label{fig:grpo2:klest}
\end{figure}

The most direct implementation, our \code{kl\_kind="seq\_simple"}, is the
simplest possible \emph{sequence-level} surrogate: $\beta\cdot\operatorname{mean}_i\big(\log\pitheta(y_i)-\log\piref(y_i)\big)$
with summed log-probabilities, i.e.\ $k_1$ summed over the sequence. It needs one extra
forward pass of a frozen model copy per batch:

\begin{lstlisting}[caption={Adding a KL term (sequence surrogate and the importance-weighted variant).},label={lst:grpo2:kl}]
ref_model = copy.deepcopy(model).eval()    # frozen anchor: +1 weight copy
for p in ref_model.parameters():
    p.requires_grad_(False)
...
with torch.no_grad():
    ref_logp = sequence_logprob(ref_model, ids, prompt_len)  # per rollout
...
pg_loss = -(adv.detach() * logps).mean()
kl_loss = kl_coeff * (logps - ref_logps).mean()      # simple surrogate
# importance-weighted (rollouts came from pi_old, not pi_theta):
w = torch.exp(logps - old_logps).detach()
kl_loss_is = kl_coeff * (w * (logps - ref_logps)).mean()
loss = pg_loss + kl_loss
\end{lstlisting}

\begin{napkinmath}[what the reference model costs]
Memory: one more bf16 copy of the weights (1.2\,GB for 0.6B, 16\,GB for 8B, 140\,GB for 70B).
Compute: one more forward pass over every rollout, $2N$ FLOPs per token, about one third of a
training step's $6N$. For large models, reference log-probs are often computed by the
\emph{rollout} workers or a dedicated reference worker (\cref{ch:rlsys}). With LoRA training
(\cref{ch:instruct}) the reference is free: it is the same weights with the adapter
disabled.
\end{napkinmath}

\subsection{Training with a KL term}\label{sec:grpo2:kltrain}

This simple surrogate can fail instructively when it is added to a long GRPO run on a small
model such as Qwen3-0.6B (with $\beta$ around $0.02$). Accuracy first improves as usual. Then the loss explodes,
rewards collapse to zero, MATH-500 accuracy follows, and inspecting the samples shows random
code-like token salad. Understanding \emph{why} is more useful than the fix:
\begin{enumerate}
\item \textbf{Wrong distribution.} The KL is an expectation under $\pitheta$, but the rollouts
were sampled from $\piold$. On those fixed samples, with $\log\piref$ a constant, minimising
$\log\pitheta(y)-\log\piref(y)$ simply means making the sampled sequences \emph{less likely}.
That is not ``stay close to the reference''. It is ``move away from whatever you just did''.
\item \textbf{No length normalisation.} Summed log-probabilities make the term scale with
response length, so long rollouts dominate it.
\item \textbf{It is the last gradient standing.} Once rewards collapse (all groups wrong),
every advantage is zero and the policy term vanishes. The KL surrogate is then the
\emph{only} source of gradient, and it keeps pushing probability away from sampled tokens,
which flattens the distributions and raises entropy. That feeds more garbage into the next
batch, a self-reinforcing collapse.
\end{enumerate}

Several remedies follow from that diagnosis: a much smaller $\beta$ (often not enough on its own); length-normalised (per-token) log-probabilities; the
unbiased $k_3$ estimator computed per token; the importance-weighted form of DeepSeek-V3.2,
$\beta\,\E_{\piold}[\rho\cdot(\log\pitheta-\log\piref)]$, which corrects for sampling from
$\piold$; or a small auxiliary reward (\cref{sec:grpo2:format}) so the policy term never
vanishes entirely. Or remove the term. For math RLVR, Dr.~GRPO, the OLMo~3 team and
DeepSeek-V3.2 all report equal or better results with the KL weight at or near zero, and
dropping it also removes the reference model's memory and compute.

\paragraph{What the toy shows.} Neither KL variant destabilised our short runs.
With $\beta=0.02$, the sequence surrogate reached greedy accuracy 0.93 and 0.98 and per-token
$k_3$ (inside the token-level clipped loss) 0.91 and 0.97, both statistically
indistinguishable from the baseline's 0.89 and 0.96. That is consistent with the diagnosis
above. The surrogate is only dangerous once rewards collapse and it becomes the sole
gradient, and on the toy task rewards rise monotonically, so mechanism (3) never fires.
Exercise~38.5 isolates the mechanism directly by training on a batch with zero advantages.
Two systems observations from these runs: the reference forward pass doubled the measured
no-grad log-prob phase (from about 0.3 to 0.6 times the update phase), roughly 13\% of a
whole step on this workload, close to the one-third-of-a-training-step estimate above; and the logged KL values differ by a factor of 4--5 between estimators
(about 2.5 nats per \emph{sequence} for the summed surrogate, 0.5--0.6 nats per \emph{token}
for $k_3$), so a $\beta$ tuned for one estimator means something very different for
another.


% ======================================================================
\section{Adding an explicit format reward}\label{sec:grpo2:format}

So far the only reward has been correctness, with the \verb|\boxed{}| requirement as an
implicit format check. That is sufficient for training reasoning models. Most recipes also add
\emph{auxiliary} rewards: a format reward, a length penalty, a language-consistency reward
(DeepSeek-R1 penalised mixing languages within a chain of thought). The canonical format
reward asks the model to put its reasoning between \verb|<think>| and \verb|</think>| and the
final answer after it.

\subsection{Thinking tokens}\label{sec:grpo2:think}

\verb|<think>| and \verb|</think>| are ordinary vocabulary entries that delimit the
reasoning. They are useful because they let a serving system hide, truncate, or budget the
reasoning separately from the answer, and because evaluation can parse the answer reliably.
They only work well as \emph{single} tokens. The Qwen3 \emph{base} tokenizer splits
\verb|<think>| into three sub-word pieces, while the Qwen3 reasoning tokenizer maps
\verb|<think>| and \verb|</think>| to the dedicated ids \num{151667} and \num{151668}. Like
most modern tokenizers, the base vocabulary reserves unused slots: its embedding matrix has
\num{151936} rows, more than the number of tokens the base tokenizer emits, so the special
tokens can be registered with the tokenizer (\code{add\_special\_tokens}) without resizing the
model. Their embeddings, however, are untrained in the base model. A base model has never seen
these ids and will not produce them reliably until some SFT or the reasoning variant's
training has taught it to. Our toy tokenizer (\code{ToyTokenizer}) has \verb|<think>| and
\verb|</think>| as single tokens from the start, and its SFT data uses them in 20\% of
demonstrations.

The reward itself is a two-line check on token ids: 1 if \verb|<think>| occurs and is
followed later by \verb|</think>|, else 0.

\begin{lstlisting}[caption={A token-level format reward and its combination with correctness.},label={lst:grpo2:format}]
def reward_format(token_ids, prompt_len, start_id, end_id):
    gen = token_ids[prompt_len:].tolist()
    try:
        return float(gen.index(start_id) < gen.index(end_id))
    except ValueError:                       # one of the tags is missing
        return 0.0

def total_reward(correct, fmt, weight=1.0, conditional=False):
    if conditional and correct <= 0:
        return correct                       # no format credit for wrong answers
    return correct + weight * fmt
\end{lstlisting}

The prompt should also ask for the format, for example: ``Write your reasoning between
\verb|<think>| and \verb|</think>|, then give the final result as \verb|\boxed{ANSWER}|.''

\begin{figure}[ht]
\centering
\resizebox{\linewidth}{!}{%
\begin{tikzpicture}[c/.style={draw=gray!60,minimum height=6.5mm,font=\scriptsize},>={Stealth[length=1.6mm]}]
\node[font=\small\bfseries] at (4.3,0.7) {rollout};
\node[font=\small\bfseries] at (9.0,0.7) {correct};
\node[font=\small\bfseries] at (10.5,0.7) {format};
\node[font=\small\bfseries,align=center] at (12.4,0.7) {sum\\{\scriptsize $w=1$}};
\node[font=\small\bfseries,align=center] at (14.4,0.7) {conditional\\{\scriptsize $w=1$}};
\foreach \y/\t/\c/\f/\s/\cs/\cc/\fc in {
 0/{<think>4+3=7</think>\textbackslash boxed\{7\}}/1/1/2/2/keygreen!25/napkin!30,
 -0.75/{\textbackslash boxed\{7\}}/1/0/1/1/keygreen!25/napkin!4,
 -1.5/{<think>4+3=8</think>\textbackslash boxed\{8\}}/0/1/1/0/keygreen!4/napkin!30,
 -2.25/{8}/0/0/0/0/keygreen!4/napkin!4} {
  \node[c,minimum width=56mm,font=\scriptsize\ttfamily] at (4.3,\y) {\t};
  \node[c,minimum width=10mm,fill=\cc] at (9.0,\y) {\c};
  \node[c,minimum width=10mm,fill=\fc] at (10.5,\y) {\f};
  \node[c,minimum width=14mm] at (12.4,\y) {\s};
  \node[c,minimum width=14mm] at (14.4,\y) {\cs};
}
\node[font=\scriptsize,warnred,align=left,anchor=west] at (15.3,-1.5) {format credit\\for a wrong answer};
\draw[->,warnred] (15.3,-1.5)--(12.95,-1.5);
\end{tikzpicture}}
\caption{Combining a correctness reward with a format reward on four rollouts of the toy
prompt. Plain summation gives the wrong-but-well-formatted third rollout the same total as
the correct-but-unformatted second one, so their advantages are equal. Conditioning the
format credit on correctness removes that incentive.}
\label{fig:grpo2:formatfig}
\end{figure}

\subsection{Training a model to emit thinking tokens}\label{sec:grpo2:formattrain}

Two practical caveats come first. You cannot reasonably RL-train a model to emit tokens it
has never produced. If the probability of \verb|<think>| is $10^{-9}$, no rollout will contain it,
the format reward is always 0 and its gradient is always zero. Format rewards therefore
\emph{reinforce} a behaviour the model already has at some rate, from SFT or from its reasoning
variant's training. Format-reward experiments are therefore best run on the
Qwen3 reasoning variant, not the base model.

Second, every auxiliary reward is a new way to be rewarded without being right. With weight 1,
a well-formatted wrong answer earns as much as an unformatted right one
(\cref{fig:grpo2:formatfig}). A run of this kind on the reasoning variant can show the failure clearly: the format
reward climbs to near-saturation while MATH-500 accuracy falls and responses get shorter: the
model learns that the tags were the cheapest reward on offer. Three fixes are standard:
\begin{itemize}
\item \textbf{Down-weight} the auxiliary reward (e.g.\ $w=0.1$) so it breaks ties without
competing with correctness.
\item \textbf{Make it conditional}: format credit only for correct answers.
\item \textbf{Normalise each component separately} (GDPO, \cref{sec:grpo2:more}): standardise
correctness and format within the group \emph{before} combining them, so a high-variance
format signal cannot drown a low-variance correctness signal in the shared advantage.
\end{itemize}
Note also an artefact visible in such runs: with one inner update and no KL, every ratio is 1
and the advantages average to zero, so the logged scalar loss is identically zero. Monitor
\code{reward\_avg}, the per-component averages, \code{adv\_std} and entropy instead.

\paragraph{What the toy shows.} Because 20\% of the SFT demonstrations already used
\verb|<think>| tags, every format variant drove \code{format\_avg} to about 0.99 within
roughly 30 steps, and responses grew from about 4.5 to 12 tokens to carry the scratchpad. The
interesting quantity is what happened to \emph{correctness}. With weight 1 the late-run
correct rate was 0.72--0.76 (greedy accuracy 0.84--0.87), below the baseline's 0.82--0.83
(0.89--0.96). Weight 0.1 was no better (0.74, greedy 0.83--0.84): even a small weight breaks
ties between otherwise equal rollouts, and on this task the tags are much easier to learn
than arithmetic. Conditioning the format credit on correctness recovered the baseline
(correct rate 0.80--0.86, greedy accuracy 0.91--0.96). GDPO's per-component normalisation
landed in between (0.75, greedy 0.87--0.91). Its logged ``reward'' is identically zero,
because the combined signal is a sum of z-scores, a reminder to log raw components alongside
whatever the optimiser sees.


\subsection{More GRPO modifications}\label{sec:grpo2:more}

We now have the original GRPO objective in full: clipped ratios, a KL term and composite
rewards. Since DeepSeek-R1, a stream of papers has modified almost every term of
\cref{eq:grpo:full}. The changes are easiest to keep straight by asking \emph{which part of the
objective} each one touches (\cref{fig:grpo2:taxonomy}). In the list below, the parentheses
name the work that proposed or popularised each change, and the code column of
\cref{tab:grpo2:mods} points to its implementation in \code{llmsys.rl}.

\begin{figure}[ht]
\centering
\resizebox{\linewidth}{!}{%
\begin{tikzpicture}[cat/.style={draw,rounded corners=3pt,font=\small\bfseries,text width=34mm,minimum height=9mm,align=center},
  it/.style={draw=gray!40,rounded corners=2pt,font=\small,align=left,text width=105mm,minimum height=9mm,anchor=west}]
\foreach \y/\c/\t/\i in {
 0/fdata/{which data get gradient}/{(1) drop zero-variance groups \quad (2) active / dynamic sampling},
 -1.15/keygreen/{advantages \& rewards}/{(7) no std scaling \quad (12) keep std scaling \quad (13) per-reward normalisation},
 -2.3/fcompute/{loss aggregation}/{(3) token-level mean instead of per-sequence mean},
 -3.45/sysblue/{ratios, clipping, IS weights}/{(5) clip-higher \quad (6) truncated IS for engine mismatch \quad (10) mask off-policy sequences\\(11) keep the sampling mask \quad (14) sequence-level ratio \quad (15) clip the IS weight},
 -4.6/gray/{KL term}/{(4) drop it \quad (8) per-domain $\beta$ \quad (9) importance-weighted KL}} {
  \node[cat,fill=\c!15,draw=\c] (c) at (0,\y) {\t};
  \node[it] at (2.2,\y) {\i};
}
\end{tikzpicture}}
\caption{Fifteen post-R1 modifications of GRPO, grouped by the part of the objective they
change. Numbers refer to the list in the text.}
\label{fig:grpo2:taxonomy}
\end{figure}

\begin{enumerate}
\item \textbf{Filter zero-gradient groups} (DAPO). Drop prompts whose $G$ rollouts all
received the same reward. They contribute nothing to the gradient but still dilute averages
and occupy the batch.
\item \textbf{Active (dynamic) sampling} (DAPO). Instead of merely dropping such prompts,
over-sample and keep drawing until the batch is full of mixed groups, so every step has a
full batch of signal.
\item \textbf{Token-level loss} (DAPO). Average the per-token loss over \emph{all tokens in
the batch} rather than first within each sequence and then over sequences. With per-sequence
averaging each token of a 2000-token rollout weighs 100 times less than a token of a 20-token
rollout, so long wrong answers are under-penalised (and long right ones under-rewarded).
\item \textbf{Drop the KL term} (DAPO, Dr.~GRPO). For verifiable rewards the anchor mostly
slows learning, and dropping it saves a model copy.
\item \textbf{Clip-higher} (DAPO). Use $\varepsilon_{\text{high}}>\varepsilon_{\text{low}}$
(e.g.\ 0.28 vs.\ 0.2) so low-probability tokens with positive advantage can grow faster,
which counteracts entropy collapse.
\item \textbf{Truncated importance sampling} (verl). Rollout engines (vLLM, SGLang) compute
probabilities with different kernels, precision and batching than the trainer, so even ``on-policy''
samples come from a slightly different distribution $\pi_{\text{rollout}}\neq\piold$. Weight
each token by $\min(\piold/\pi_{\text{rollout}},C)$ to correct for it without letting rare
tokens explode the update.
\item \textbf{No std normalisation of advantages} (Dr.~GRPO). Dividing by the group std
up-weights prompts that are nearly always solved or nearly always failed (small std), a
difficulty bias. Dr.~GRPO uses $A_i=r_i-\bar r$ and also replaces the per-sequence length
normalisation by a constant, removing a length bias.
\item \textbf{Per-domain KL strength} (DeepSeek-V3.2). Keep the KL for domains where drift is
harmful and set $\beta=0$ for math.
\item \textbf{Importance-weighted KL} (DeepSeek-V3.2). Multiply the KL estimate by
$\pitheta/\piold$ because the samples come from $\piold$. This is exactly the bias diagnosed in
\cref{sec:grpo2:kltrain}.
\item \textbf{Mask badly off-policy sequences} (DeepSeek-V3.2). Exclude negative-advantage
sequences whose average $\log\piold-\log\pitheta$ exceeds a threshold. Pushing down
samples the current policy would hardly generate anyway is mostly noise.
\item \textbf{Keep the sampling mask} (DeepSeek-V3.2). If rollouts used top-$p$/top-$k$,
compute $\piold$ and $\pitheta$ over the \emph{same truncated support}, so the ratio compares
like with like.
\item \textbf{Keep the original advantage normalisation} (DeepSeek-V3.2). A counterpoint to
(7): standard group-normalised advantages worked well in their setting, so ablate before
removing.
\item \textbf{Normalise each reward component before aggregation} (GDPO). Standardise each
reward type within the group, then combine, so components with different scales and variances
contribute comparably.
\item \textbf{Sequence-level ratio and clipping} (GSPO). Use the length-normalised sequence
ratio $s_i=\exp\big(\frac1{|y_i|}\sum_t\log\rho_{i,t}\big)$ with a tight clip (around
$3$--$4\times10^{-4}$). This matches the sequence-level reward and is reported to be more
stable, notably for MoE models where token-level ratios are noisy.
\item \textbf{Clip importance weights, not updates} (CISPO, MiniMax-M1). Use
$-\operatorname{sg}(\min(\rho_{i,t},1+\varepsilon))\,A_i\log\pitheta(y_{i,t})$: every token
keeps a gradient (PPO-clip silences clipped tokens entirely, which can drop rare but
important ``reflection'' tokens), but its weight is bounded.
\end{enumerate}

\begin{table}[ht]
\centering\small
\begin{tabularx}{\linewidth}{@{}lXl@{}}
\toprule
\# & One-line summary & \code{GRPOConfig} / function\\
\midrule
1, 2 & spend gradient only on mixed groups & \code{filter\_zero\_std}\\
3 & mean over all tokens, not per sequence & \code{agg="token\_mean"}\\
4, 8 & KL optional / per domain & \code{kl\_coeff=0}\\
5 & asymmetric clip & \code{eps\_low}, \code{eps\_high}\\
6 & engine/trainer mismatch weights & \code{truncated\_is\_weights}\\
7 & centre but don't scale & \code{adv\_scale="none"}, \code{agg="seq\_mean\_token\_sum\_norm"}\\
9 & IS-weighted KL & \code{kl\_is\_weighted=True}\\
10 & drop very off-policy negatives & \code{offpolicy\_delta}\\
11 & same support for both policies & \code{top\_p\_renormalized\_logprob}\\
13 & per-component normalisation & \code{gdpo=True}\\
14 & geometric-mean sequence ratio & \code{loss="gspo"}\\
15 & bounded, gradient-stopped IS weight & \code{loss="cispo"}\\
\bottomrule
\end{tabularx}
\caption{Where each modification lives in the companion code.}
\label{tab:grpo2:mods}
\end{table}

\begin{lstlisting}[caption={Token-level clipped GRPO loss with the main options (abridged from \texttt{llmsys/rl/losses.py}).},label={lst:grpo2:token}]
def token_clipped_loss(lp_new, lp_old, adv, mask, eps_low=0.2, eps_high=0.2,
                       lp_ref=None, kl_coeff=0.0, kl_kind="k3",
                       kl_is_weighted=False, agg="seq_mean_token_mean"):
    ratio = torch.exp((lp_new - lp_old.detach()) * mask)        # (N, T)
    a = adv.detach()[:, None]
    obj = torch.minimum(ratio * a,
                        torch.clamp(ratio, 1 - eps_low, 1 + eps_high) * a)
    per_tok = -obj
    if lp_ref is not None and kl_coeff > 0:
        kl = kl_estimate(lp_new, lp_ref, kl_kind)               # k1 / k2 / k3
        if kl_is_weighted:                                      # DeepSeek-V3.2
            kl = kl * ratio.detach()
        per_tok = per_tok + kl_coeff * kl
    return aggregate(per_tok, mask, agg)    # GRPO / DAPO / Dr. GRPO reductions
\end{lstlisting}

\paragraph{What the toy shows.} \Cref{fig:grpo2:ablation} collects every variant. The four
descendants we ran with the trainer's defaults (DAPO-style token-level loss with clip-higher
and zero-variance filtering, Dr.~GRPO, GSPO and CISPO) all landed within noise of the
baseline: greedy accuracy 0.91--0.96 against 0.93. That is the expected outcome, not a
disappointment. These modifications target failure modes that a 76k-parameter model on
five-token answers never reaches: long-sequence length bias, entropy collapse over thousands
of steps, MoE routing noise, and stale samplers. Their value shows up in long runs on real
models, which is where to evaluate them (Lab~38's stretch goal). The toy's job here is to
verify that each implementation is \emph{neutral when it should be}. A modification that hurt
on a problem this easy would be a bug.

\begin{figure}[ht]
\centering
\begin{tikzpicture}
\begin{axis}[xbar,width=0.82\linewidth,height=10.5cm,bar width=4pt,xmin=0,xmax=1.05,
  ytick={14,13,12,11,10,9,8,7,6,5,4,3,2,1,0},yticklabels={{baseline (Ch.~27)},{4 passes, no clip},{4 passes, clip},{8 steps, lr$\times$3, no clip},{8 steps, lr$\times$3, clip},{KL: sequence surrogate},{KL: per-token $k_3$},{format $w{=}1$},{format $w{=}0.1$},{format, conditional},{format, GDPO},{DAPO-style},{Dr.~GRPO},{GSPO},{CISPO}},y tick label style={font=\scriptsize},
  x tick label style={font=\scriptsize},xlabel={score},enlarge y limits=0.04,
  legend style={font=\scriptsize,at={(0.5,1.01)},anchor=south,legend columns=2,draw=none},
  xmajorgrids,grid style={gray!20}]
\addplot[fill=sysblue!60,draw=none,error bars/.cd,x dir=both,x explicit] coordinates {(0.925,14) +- (0.035,0) (0.560,13) +- (0.060,0) (0.915,12) +- (0.035,0) (0.085,11) +- (0.015,0) (0.145,10) +- (0.015,0) (0.955,9) +- (0.025,0) (0.940,8) +- (0.030,0) (0.855,7) +- (0.015,0) (0.835,6) +- (0.005,0) (0.935,5) +- (0.025,0) (0.890,4) +- (0.020,0) (0.910,3) +- (0.040,0) (0.950,2) +- (0.020,0) (0.930,1) +- (0.040,0) (0.955,0) +- (0.015,0)};
\addplot[fill=keygreen!55,draw=none,error bars/.cd,x dir=both,x explicit] coordinates {(0.826,14) +- (0.007,0) (0.550,13) +- (0.015,0) (0.803,12) +- (0.008,0) (0.087,11) +- (0.029,0) (0.145,10) +- (0.007,0) (0.855,9) +- (0.002,0) (0.847,8) +- (0.015,0) (0.739,7) +- (0.020,0) (0.743,6) +- (0.001,0) (0.827,5) +- (0.029,0) (0.747,4) +- (0.001,0) (0.825,3) +- (0.005,0) (0.853,2) +- (0.005,0) (0.846,1) +- (0.001,0) (0.838,0) +- (0.030,0)};
\legend{greedy accuracy after 150 steps,late-run train correct rate}
\draw[gray,dashed] (axis cs:0.925,-1) -- (axis cs:0.925,15);
\end{axis}
\end{tikzpicture}
\caption{Every GRPO variant of this chapter on the toy task: mean over two seeds, with the bars'
whiskers spanning the two seeds; 150 steps from the same SFT start; measured on CPU. Dashed
line: baseline mean greedy accuracy. With only two seeds, differences under about 0.05 are noise.
Clear effects: unclipped rollout reuse hurts, clipping fixes moderate reuse but not an
excessive step size, and an unconditional format reward costs correctness while a conditional
one does not.}
\label{fig:grpo2:ablation}
\end{figure}

% ======================================================================
\section{The systems cost of each knob}\label{sec:grpo2:costs}

Algorithm papers report accuracy. Whoever runs the job also has to pay for it.
\Cref{tab:grpo2:costs} summarises how each mechanism of this chapter changes memory, compute
and scheduling. Most of the changes have a positive effect.

\begin{table}[ht]
\centering\small
\begin{tabularx}{\linewidth}{@{}lX@{}}
\toprule
Mechanism & Cost / effect on the system\\
\midrule
entropy metric & full-vocabulary reduction; compute in chunks on the no-grad pass, otherwise
$(NT)\times|\mathcal V|$ fp32 memory\\
clipping + rollout reuse & \emph{saves} sampling: $k$ updates per batch divide sampling cost per
update by $k$; enables asynchronous samplers\\
KL to reference & +1 model copy in memory, +1 forward ($\approx$\,33\% of a training step);
free with LoRA (disable adapter)\\
dropping KL & removes both of the above\\
format/aux rewards & negligible (string or id checks on CPU); code-execution rewards need a
sandbox fleet\\
zero-variance filtering & wasted rollouts: a fraction $p^G+(1-p)^G$ of sampling is discarded\\
dynamic sampling & over-samples by $\approx1/(1-\text{waste})$ and makes step time
\emph{variable}, which complicates synchronous pipelines\\
token-level aggregation & none (a different reduction); changes effective learning rate per
token, so retune lr\\
truncated IS (engine mismatch) & needs the sampler to return per-token log-probs (cheap) and
the trainer to recompute them (already done)\\
sampling-mask preservation & needs the top-$p$ support per token, either stored (huge) or
recomputed from the trainer's own logits (approximate)\\
GSPO / CISPO & none beyond the base loss; GSPO's sequence ratio needs all tokens of a sequence
in one micro-batch\\
\bottomrule
\end{tabularx}
\caption{What the stabilisation mechanisms cost. ``Saves'' entries are why clipping-based
off-policy training underlies essentially all high-throughput RL systems (\cref{ch:rlsys}).}
\label{tab:grpo2:costs}
\end{table}

\begin{napkinmath}[the price of dynamic sampling]
Suppose half of a math dataset's prompts are solved with $p\approx0.9$ by the current policy
and half with $p\approx0.3$. With $G=16$, zero-variance probabilities are
$0.9^{16}+0.1^{16}\approx0.185$ and $0.3^{16}+0.7^{16}\approx0.003$, so about 9.4\% of groups
are wasted. To fill a batch of 512 mixed groups, dynamic sampling must generate about
$512/0.906\approx565$ groups, or $\approx\num{9000}$ rollouts. As the policy improves, the easy
half drifts to $p\approx0.97$ ($0.97^{16}\approx0.61$) and the waste rises to about 30\%, so the
over-sampling factor grows to 1.45. Curricula that retire mastered prompts keep this
factor near 1.
\end{napkinmath}

% ======================================================================
\section{Summary}\label{sec:grpo2:summary}
\begin{itemize}
\item Read a GRPO run through several metrics together: moving-average reward, response length,
held-out accuracy, advantage mean (sanity) and std (signal), zero-variance fraction, entropy,
ratio and clip fraction. The loss is mainly a sanity check.
\item Reward stuck at 1 or 0 means no signal. Entropy collapse means no exploration. Entropy
explosion means a damaged model.
\item Importance ratios correct for training on samples from an older snapshot, and clipping
keeps each batch's updates inside a trust region. Together they make rollout reuse and
asynchronous sampling safe, which is a throughput win.
\item The KL penalty anchors the run to a reference policy. A naive sampled surrogate can
destabilise training once rewards collapse. Per-token $k_3$, importance weighting, or simply
$\beta=0$ (common for math) avoid it.
\item Auxiliary rewards such as \verb|<think>| formatting are cheap but hackable. Down-weight
them, condition them on correctness, or normalise components separately (GDPO).
\item Post-R1 variants each change one term: what data gets gradient (DAPO filtering and dynamic
sampling), advantage scaling (Dr.~GRPO), loss aggregation (token-level), ratio and clipping
(clip-higher, GSPO, CISPO, truncated IS), and KL handling (DeepSeek-V3.2).
\item Every knob has a systems cost. The cheapest stabiliser, clipping, is also the one that
unlocks most of the throughput.
\end{itemize}

% ======================================================================
\section{Exercises}\label{sec:grpo2:exercises}

\begin{exercise}[testing the format reward]
Run \code{reward\_format} on a completion with both tags in order and confirm it returns 1.
Then introduce a typo in one tag, reverse their order, and delete one tag; each should give
0. Which of these edits would a \emph{string}-based check (\code{think\_format\_ok}) also catch?
Which tokenisation edge case can break the id-based check but not the string-based one?
\end{exercise}

\begin{exercise}[conditional format reward]
Run the four format variants (\code{format\_w1}, \code{format\_w0.1}, \code{format\_cond},
\code{format\_gdpo}) with seeds 0, 1 and 2. Report final greedy accuracy, late-run \code{correct\_avg} and \code{format\_avg}. Does
conditioning change which behaviour the policy adopts first? Explain the result with the
advantages of \cref{fig:grpo2:formatfig}.
\end{exercise}

\begin{exercise}[the clip fraction]
For the \code{clip\_reuse4} variant, plot \code{clip\_frac} and \code{policy\_ratio} over
training. Vary \code{inner\_epochs} from 1 to 8 and find the value where the clip fraction
exceeds 0.2. Relate this to the wall-clock saving of rollout reuse.
\end{exercise}

\begin{exercise}[unbiasedness of $k_3$]
Show that $\E_{y\sim p}[k_3]=\KL(p\,\|\,q)$ where $k_3=r-1-\log r$ and $r=q(y)/p(y)$, and
that $k_3\ge0$ pointwise. Verify numerically for a random pair of categoricals over 1000
outcomes, and compare the empirical variance of $k_1$, $k_2$ and $k_3$ with 64 samples.
\end{exercise}

\begin{exercise}[diagnose the KL collapse]
On a fixed batch where all advantages are zero, take 50 optimiser steps on the
sequence-level KL surrogate alone. Plot the mean entropy and the probability of the sampled
tokens. Repeat with per-token $k_3$. Explain the difference in terms of what each loss's
gradient does to the sampled tokens.
\end{exercise}

\begin{exercise}[aggregation and length]
Construct a group of two rollouts with $A=(+1,-1)$ and lengths 10 and 1000. Compute, under
each of the three aggregations of \code{aggregate\_token\_loss}, the total gradient weight
received by each rollout and by each of its tokens. Which rule lets a long wrong answer off
lightly?
\end{exercise}

\begin{exercise}[Dr.~GRPO's difficulty bias]
For binary rewards and $G=8$, compute the gradient-weight multiplier
$\sum_i|A_i|$ as a function of $p$ with and without std scaling. Which prompts does std
scaling up-weight? Argue whether that is a feature (focus on edge cases) or a bug.
\end{exercise}

\begin{exercise}[chunked entropy]
Measure peak memory of \code{token\_entropy} on a $(4096, \num{151936})$ bf16 logits tensor
versus \code{chunked\_logprob\_entropy} with chunk sizes 256 and 1024 (use
\code{torch.cuda.}\allowbreak\code{max\_memory\_allocated} or \code{tracemalloc}). Plot memory and time against
chunk size.
\end{exercise}

\begin{exercise}[napkin: reference-model placement]
For a 32B policy trained with FSDP on 16 GPUs, compare three placements of the KL
reference: (a) a full bf16 copy on every trainer GPU, (b) FSDP-sharded alongside the policy,
(c) computed by the rollout workers. Estimate memory per GPU and extra communication per step
for each.
\end{exercise}

\begin{exercise}[engine mismatch]
Simulate an inference engine by sampling rollouts from the toy model with its weights cast to
bf16 while the trainer uses fp32. Measure the distribution of $\log\piold-\log\pi_{\text{rollout}}$
per token. Train with and without \code{truncated\_is\_weights} and compare stability.
\end{exercise}

% ======================================================================
\section{Lab 38: making GRPO robust}\label{sec:grpo2:lab}

\begin{labbox}{38}{Diagnostics, clipping, KL, format rewards and GRPO's descendants}
\labdir{lab38_grpo_plus}
\labsection{Goals}
Implement the diagnostics and stabilisers of this chapter as small, testable functions, then
use the library trainer to ablate them on the toy task and interpret the results like a
practitioner.

\labsection{Setup}
\begin{lstlisting}[style=shell]
cd code
LAB_IMPL=starter python -m pytest -q labs/lab38_grpo_plus
python labs/lab38_grpo_plus/run.py --list
\end{lstlisting}

\labsection{Tasks} (functions in \code{starter.py})
\begin{enumerate}
\item Diagnostics: \code{token\_entropy}, \code{sequence\_logprob\_and\_entropy},
\code{chunked\_logprob\_entropy}, \code{advantage\_stats}.
\item Clipping: \code{policy\_ratios}, \code{clipped\_policy\_loss} (returning the clip
fraction). The tests check both the value and \emph{which} samples receive gradient.
\item KL: \code{kl\_k1}, \code{kl\_k2}, \code{kl\_k3}, \code{kl\_sequence\_surrogate},
\code{kl\_importance\_weighted}.
\item Rewards: \code{reward\_format}, \code{total\_reward} (weighted and conditional).
\item Descendants: \code{dr\_grpo\_advantages}, \code{aggregate\_token\_loss} (three
reductions), \code{zero\_variance\_keep}, \code{expected\_wasted\_fraction}.
\end{enumerate}

\labsection{Measurements}
\begin{enumerate}
\item Run all variants with three seeds (about 30\,s per run on one CPU core):
pass every name from \code{run.py --list} to \code{--variants} with \code{--seeds 0,1,2}. Reproduce \cref{fig:grpo2:ablation} with
error bars.
\item For \code{baseline}, \code{noclip\_reuse4} and \code{clip\_reuse4}, plot reward, entropy
and zero-variance fraction against step and against \emph{wall-clock time}. Rollout reuse
should win on the second axis if it is stable.
\item Record \code{t\_rollout}, \code{t\_logprob}, \code{t\_update} per variant and attribute
the overhead of the KL variants to the reference forward pass.
\end{enumerate}

\labsection{Deliverables}
Passing tests; an ablation table (mean $\pm$ std over seeds); a half-page ``run review'' that
diagnoses one failed configuration from its metrics alone, before looking at its samples.

\labsection{Stretch goals}
\begin{itemize}
\item Implement DAPO's \emph{active} sampling (keep sampling until the batch has $P$ mixed groups)
and measure the over-sampling factor over training (compare with the napkin math above).
\item Add DAPO's soft overlong penalty (\code{llmsys.rl.rewards.overlong\_penalty}) and a
length-capped toy variant that requires long \verb|<think>| traces.
\item Port the best configuration to Qwen3-0.6B on MATH with Lab~37's \code{run.py qwen3}
path and evaluate on MATH-500.
\end{itemize}
\end{labbox}
